% !TeX root = ../Thesis.tex

%************************************************
\chapter{Einleitung}\label{ch:introduction}
%************************************************
\glsresetall % Resets all acronyms to not used

%**********Motivation und Kontext*************
\section{Motivation und Kontext}
\label{sec:firstsection}
Antike Münzen wurden mit handgefertigten Stempeln geprägt (vgl. \cref{sec:glnumismatik}). Münzen desselben Stempels sollten sich daher besonders ähnlich sein, und aus den Kombinationen der Stempel lässt sich die Reihenfolge der Prägung rekonstruieren \cite{wikipedia-numismatik}. Manuelle Stempelstudien sind jedoch so aufwendig, dass große Prägungen, etwa die des Römischen Reichs, kaum umfassend untersucht werden können \cite{heinecke2021unsupervised}. \todo{Rückfrage: Soll hier ein Beleg stehen, dass Stempelvergleiche eine etablierte Methode sind (z. B. Kluge, Howgego)?} Digitale Datenbanken wie das Corpus Nummorum \cite{peter2024corpus} stellen große Mengen an Münzbildern bereit (vgl. \cref{sec:gesamtansatzdatengrundlage}). Bildbasierte Verfahren wurden für antike Münzen bereits eingesetzt, etwa zur Zuordnung von Münzen zu ihrer Prägung (engl. Issue) \cite{guo2023siamese}, zur Erkennung einzelner Motive wie Adler oder Schild \cite{cooper2020learning}, zur Zeichenerkennung auf Legenden \cite{anwar2026reading} und zur Automatisierung der Stempelanalyse \cite{heinecke2021unsupervised}. Mehrere dieser Verfahren erfordern ein Training auf münzspezifischen Daten \cite{guo2023siamese,cooper2020learning,anwar2026reading}.

Die vorliegende Arbeit untersucht, ob sich einzelne Merkmale auf antiken Münzbildern auch ohne münzspezifisches Training detektieren und für einen Vergleich nutzen lassen. Entscheidend ist dabei nicht, ob ein detektiertes Merkmal inhaltlich korrekt benannt wird, sondern ob Merkmale derselben Klasse, die auf verschiedenen Münzen detektiert werden, miteinander verglichen werden können. Eine ausführliche Einordnung in bestehende Arbeiten folgt in \cref{sec:bestehendearbeiten}.


%**********Zielsetzung und Forschungsfragen*************
\section{Zielsetzung und Forschungsfragen}
\label{sec:zielsetzungundforschungsfragen}
Ziel der vorliegenden Arbeit ist es, visuelle Merkmale auf antiken Münzbildern mithilfe eines YOLO-Modells (zu erkennen und) zu lokalisieren und anschließend anhand ihrer extrahierten Feature-Vektoren miteinander zu vergleichen.
Im Rahmen der Recherche wurde kein Modell gefunden, das für die Detektion der hier betrachteten Merkmale auf antiken Münzbildern trainiert wurde \todo{Formulierung prüfen, siehe Anmerkung} (vgl. \cref{sec:modelle}). Ein eigenes Training würde zudem Annotationen der Merkmalspositionen voraussetzen, die nicht vorliegen. Deshalb wird bewusst ein Ansatz ohne Nachtraining verfolgt. Untersucht wird dabei, inwieweit sich ein nicht auf Münzen angepasstes Modell für die Detektion und den Vergleich von Merkmalen auf antiken Münzbildern eignet.
Aus dieser Zielsetzung ergeben sich folgende Forschungsfragen:
\begin{enumerate}
    \item Inwieweit lassen sich visuelle Merkmale auf antiken Münzbildern mithilfe eines nicht speziell trainierten YOLO-Modells detektieren und lokalisieren?
    \item Inwiefern lassen sich lokalisierte Merkmale anhand ihrer extrahierten Feature-Vektoren sinnvoll miteinander vergleichen?
    \item Welche Grenzen zeigt ein Ansatz ohne münzspezifisches Training bei der Anwendung auf antike Münzbilder?
\end{enumerate}

\noindent Nicht Gegenstand der Arbeit sind das Training oder Nachtrainieren eigener Modelle sowie die numismatische Bestimmung oder Datierung von Münzen. Der Vergleich ganzer Münzbilder ist in der Anwendung als ergänzende Funktion umgesetzt, wird aber nicht untersucht (vgl. \cref{sec:gesamtansatzdatengrundlage}).